\section{TD2 Key Techniques}

This TD is about the bridge between elementary discrete-time martingale
calculus and no-arbitrage pricing in a finite trading grid.  The notation is
the one of Chapter 1 and Appendix A of the course poly, especially
\polyref{Def.~1.1--1.4}, \polyref{Thm.~1.1}, \polyref{Prop.~1.7}, and
\polyref{Def.~A.6}.

\begin{techniques}
  \item \textbf{Conditional expectation.}  If \(X\) is integrable and
  \(Y\) is \(\F_n\)-measurable, then
  \[
    \E[YX\mid \F_n]=Y\E[X\mid \F_n]
  \]
  whenever the products are integrable.  If \(X\) is independent of
  \(\F_n\), then \(\E[X\mid \F_n]=\E[X]\).  These two rules are the main
  tools for random walks, because the new increment is independent of the
  past while the current value is measurable with respect to the past.

  \item \textbf{Martingales.}  A process \(M\) is an \(\F\)-martingale when
  it is adapted, integrable, and satisfies
  \[
    \E[M_{n+1}\mid \F_n]=M_n .
  \]
  For a centered random walk \(X_n=\sum_{k=1}^n \xi_k\), the martingale
  property comes from
  \(\E[\xi_{n+1}\mid \F_n]=0\).  Quadratic martingales are obtained by
  compensating the predictable drift:
  \[
    X_n^2-\sum_{k=1}^n \E[\xi_k^2]
  \]
  is an \(\F\)-martingale.

  \item \textbf{Martingale transforms.}  If \(M\) is a martingale and
  \(\Delta_n\) is bounded and adapted, then
  \[
    X_n=\sum_{k=1}^n \Delta_{k-1}(M_k-M_{k-1})
  \]
  is again a martingale.  This is the discrete stochastic-integral rule:
  \(\Delta_{k-1}\) is known before the increment \(M_k-M_{k-1}\), so the
  conditional mean of each new gain is zero.

  \item \textbf{AOA.}  In the course convention, an arbitrage opportunity is
  a self-financing strategy with zero initial capital, nonnegative terminal
  discounted wealth, and a positive probability of strict gain:
  
  \(V_0=0\) and \(\Vtilde_T\ge 0\) with \(\P(\Vtilde_T>0)>0\), where 
  
  see \polyref{Def.~1.3}.  Equivalent definitions allowing \(V_0\le 0\) are
  reduced to this one by adding or subtracting a pure risk-free holding.

  \item \textbf{Self-financing strategies.}  A strategy is self-financing
  when changes in wealth come only from price moves:
  \[
    \Vtilde_{t_k}-\Vtilde_{t_{k-1}}
      =\inner{\theta_{t_k}}{\Stilde_{t_k}-\Stilde_{t_{k-1}}}.
  \]
  Equivalently, rebalancing at \(t_{k-1}\) costs no money in discounted
  units:
  \[
    \inner{\theta_{t_{k-1}}}{\Stilde_{t_{k-1}}}
      =\inner{\theta_{t_k}}{\Stilde_{t_{k-1}}}.
  \]
  This is \polyref{Rem.~1.1}.

  \item \textbf{Forwards with deterministic dividends.}  A deterministic
  dividend is removed from the spot price by subtracting its present value.
  Thus the time-0 forward delivery price for maturity \(T\) is
  \[
    F_0=\left(S_0-\sum_j D_j e^{-r\bar t_j}\right)e^{rT}
  \]
  when the dividend dates \(\bar t_j\) are before \(T\).

  \item \textbf{Discrete pricing by backward induction.}  If an equivalent
  probability \(\Q\) makes \(\Stilde\) a martingale, then AOA holds by
  \polyref{Thm.~1.1}.  If a payoff \(G\) is replicable, its price is
  \[
    p^G_{t_k}
      =S^0_{t_k}\E_\Q\left[\frac{G}{S^0_T}\,\middle|\,\F_{t_k}\right],
  \]
  as in \polyref{Prop.~1.7}.  In exercises, it is usually more practical to
  start from the terminal value
  \[
    p^G_T=G
  \]
  and move one step backward.  At a node
  \((t_{k-1},S^1_{t_{k-1}}=s)\) of a binomial tree, the two next stock values
  at \(t_k\) are \(su\) and \(sd\).  With
  \[
    q=\frac{e^{r\Delta t}-d}{u-d},
  \]
  \polyref{Thm.~1.2--1.3} gives the one-step recursion
  \[
    p^G_{t_{k-1}}(s)
      =e^{-r\Delta t}\left(
        q\,p^G_{t_k}(su)
        +(1-q)\,p^G_{t_k}(sd)
      \right).
  \]
  The risky-asset holding \(\theta^1_{t_k}\), held over
  \((t_{k-1},t_k]\), is the slope between the two successor values:
  \[
    \theta^1_{t_k}(s)
      =
      \frac{p^G_{t_k}(su)-p^G_{t_k}(sd)}
           {su-sd}.
  \]
  In words: first fill the terminal payoffs, then compute prices one column
  backward, and at each node compute \(\theta^1\) from the two successor
  prices and stock values.
\end{techniques}
